\clearpage
\SourceLeaf{398}
\fancyhead[L]{\small ZUSÄTZE ZUR GEOMETRIE DER STELLUNG.}
\pdfbookmark[1]{Zusätze zur Geometrie der Stellung von Carnot. 1810.}{carnot-zusaetze}
\begin{center}
{\Large ZUSÄTZE.}\par\bigskip
ZUR GEOMETRIE DER STELLUNG VON CARNOT\par\bigskip
ÜBERSETZT VON SCHUMACHER.\quad 1810.\par\bigskip
\rule{.1\textwidth}{.3pt}
\end{center}
\par\medskip\Needspace{5\baselineskip}
\pdfbookmark[2]{I.}{carnot-zusatz-i}
\centerline{\textbf{I.}}\nobreak\smallskip
\{Folgende analytische Behandlung der merkwürdigen Punkte eines Dreiecks verdanke ich der Güte des Herrn Professor \textsc{Gauss}. \textsc{Schumacher}.\}

Es seien \(A,A',A''\) die drei Winkelpunkte eines Dreiecks und deren Coordinaten respective
\[
\begin{array}{ll}
x,&y\\x',&y'\\x'',&y''
\end{array}
\]
\par
Die Coordinaten der Punkte \(B,B',B''\), welche die Seiten \(A'A'',A''A,AA'\) halbiren, werden offenbar sein
\[
\begin{array}{ll}
\tfrac12(x'+x''),&\tfrac12(y'+y'')\\
\tfrac12(x''+x),&\tfrac12(y''+y)\\
\tfrac12(x+x'),&\tfrac12(y+y')
\end{array}
\]
\par
Man nehme auf den Linien \(AB,A'B',A''B''\) (vorwärts oder rückwärts verlängert, wenn es nöthig ist), von \(A,A',A''\) ab gezählt, Stücke, welche jenen respective proportional sind, und sich dazu wie \(n:1\) verhalten. Falls man die Stücke rückwärts nimmt, hat man \(n\) als negativ anzusehen. Dieser Stücke Endpunkte heissen \(C,C',C''\), so sind ihre Coordinaten
